\begin{tikzpicture}[node distance=0.45cm]
\node[fato] (F) at (0,0) {\textbf{Tabela de Fatos Ocupação de Turma}
\nodepart{two}(FK) chave ano letivo\\(FK) chave escola\\(FK) chave território\\(FK) chave série\\(FK) chave turno\\id turma \hfill\textit{(degenerada)}\\quantidade turmas \hfill\textit{(artificial)}\\matrículas no ano\\matrículas ativas em junho\\turmas ativas em junho\\\textit{alunos por turma (derivado)}};
\node[dim, anchor=north east] (A) at (-3.35,3.25) {\textbf{Dimensão Ano Letivo}\nodepart{two}(PK) chave ano letivo\\ano\\qualidade código escola\\observação};
\node[dim, below=of A] (E) {\textbf{Dimensão Escola}\nodepart{two}(PK) chave escola\\código escola, versão\\nome escola, tipo escola, rede\\distrito, DRE, situação\\...};
\node[dim, below=of E] (T) {\textbf{Dimensão Território}\nodepart{two}(PK) chave território\\distrito\\DRE, participação DRE\\município, UF};
\node[dim, anchor=north west] (S) at (3.35,2.35) {\textbf{Dimensão Série}\nodepart{two}(PK) chave série\\modalidade, segmento, série\\ciclo, educação especial\\código etapa, etapa, nível};
\node[dim, below=of S] (TU) {\textbf{Dimensão Turno}\nodepart{two}(PK) chave turno\\turno\\jornada\\origem turno};
\draw[liga] (A.east) -- \ladoesq{F}{1}{3};
\draw[liga] (E.east) -- \ladoesq{F}{2}{3};
\draw[liga] (T.east) -- \ladoesq{F}{3}{3};
\draw[liga] (S.west) -- \ladodir{F}{1}{2};
\draw[liga] (TU.west) -- \ladodir{F}{2}{2};
\end{tikzpicture}
